\begin{abstract}
    Traffic-related air pollution remains an important public health concern in 
    dense urban environments, but traditional monitoring infrastructure cannot 
    easily provide continuous, street-level measurements by vehicle type. This 
    thesis investigates whether existing New York City traffic-camera feeds can 
    be used to estimate traffic flow and vehicle composition for downstream emissions
    analysis. To address this problem, a 1,200-image low-resolution vehicle detection 
    dataset was constructed from 511NY and NYSDOT imagery and annotated using a 
    6-class taxonomy aligned with broad EPA MOVES vehicle categories. Several 
    object-detection architectures were evaluated, including YOLO, Faster R-CNN, 
    RT-DETR, and Deformable DETR, with YOLOv11s providing the strongest overall 
    performance. The baseline YOLOv11s model achieved an AP50 of 0.804 and an mAP 
    of 0.615 on the validation set, while the final refined model improved to an 
    AP50 of 0.857 and an mAP of 0.661 through targeted data refinement. On a 
    held-out test set, the final model reached an AP50 of 0.772 and an mAP of
    0.607, while performance on a cross-camera benchmark of unseen locations dropped
    to an AP50 of 0.658 and an mAP of 0.488, showing that generalization to new 
    camera views remained a challenge. Comparison against New York City Automated 
    Traffic Recorder (ATR) data showed that the pipeline was better at recovering 
    the temporal pattern of traffic flow than exact count magnitude, with the 
    strongest agreement at Atlantic Avenue at Boerum Place, where Pearson correlation
    reached 0.844 and MAPE was 22.96\%. Overall, the results show that public 
    traffic-camera imagery can support practical traffic-pattern and fleet-composition
    analysis, while also highlighting the need for broader training data, temporal 
    tracking, and additional operational inputs before such a system can support 
    fully calibrated traffic counting and emissions modeling.
\end{abstract}

